\ficha{Villela (2001), Debates monetários no Brasil do século XIX}{villela2001}

\begin{identificacao}
VILLELA, André. The Quest for Gold: Monetary Debates in Nineteenth-century
Brazil. \emph{Brazilian Journal of Political Economy} (Revista de Economia
Política), v. 21, n. 4 (84), p. 690-704, out.-dez. 2001.

Artigo de história do pensamento econômico e de história monetária, 15 páginas,
em inglês, com resumo em português. Lido integralmente na versão publicada.
DOI 10.1590/0101-31572001-1222.
\end{identificacao}

\bloco{Problema e tese}
O artigo pergunta o que de fato separava metalistas e papelistas no debate
monetário brasileiro de meados do século XIX. Responde que a historiografia
classificou os agentes sem perguntar quem defendia ou não a conversibilidade do
mil-réis, e que essa omissão produziu leituras erradas da ideologia da gestão
monetária do período. Aplicado o critério, quase todos os participantes,
metalistas e papelistas, defendiam a conversibilidade e o padrão-ouro. A
oposição real corria por outros eixos, a pluralidade do direito de emissão, as
regras prudenciais impostas aos bancos e a licença para expandir em pânico.
Defensores de circulação puramente fiduciária foram exceção numérica no Império.

\bloco{Argumento}
\begin{enumerate}[leftmargin=*,nosep]
\item Villela recompõe primeiro as controvérsias inglesas. A polêmica bulionista,
aberta pelo Bullion Committee de 1810, foi resolvida em 1819 a favor dos
bulionistas, com o ato que determinou a retomada da conversibilidade das notas do
Banco da Inglaterra (p. 692). Fixado o princípio, a disputa seguinte, entre
currency school (McCulloch, Loyd Overstone) e banking school (Mill, Tooke,
Fullarton), passou a ser sobre regra contra discrição, com o princípio de
currency de um lado e as doutrinas das real bills, da lei do refluxo e das
necessidades do comércio do outro (p. 692-693).

\item O passo decisivo é notar o que as duas escolas tinham em comum.

\chave{there was a crucial common thread to both schools of monetary thought. It
was the presupposition of convertibility of notes}{p. 693}

A conversibilidade em espécie à vista, e o câmbio fixo que ela implicava, eram a
ortodoxia da época, contra a qual poucos ousavam se colocar (p. 693). Daí os dois
eixos de conflito que o autor isola: quem pode emitir, isto é, direitos de
emissão, e com que lastro, isto é, conversibilidade (p. 695). O artigo trata do
segundo.

\item No Brasil, a Lei 401, de 11 de setembro de 1846, de iniciativa de Bernardo
de Vasconcelos no Senado, fixou o par legal de 27 pence por mil-réis (p. 696,
nota 13). Era paridade nocional, porque a circulação se compunha de notas do
Tesouro e de vales bancários inconversíveis. A conversibilidade veio com o
terceiro Banco do Brasil, em 1854, e vigorou por menos de seis anos, em três
intervalos entre abril de 1854 e a crise Souto de setembro de 1864 (p. 696,
Tabela 1, p. 697). Nesses intervalos as notas eram lastreadas parte em ouro,
parte em papel do Tesouro, este inconversível, de modo que o Banco do Brasil
tecnicamente convertia sua moeda em papel. Ainda assim o câmbio se manteve junto
do par, resultado que Villela atribui à ``conversibilidade potencial'', a simples
promessa de resgate, que teria construído credibilidade no sistema (p. 697).

\item A importação das categorias é atribuída à historiografia, não aos
contemporâneos: metalistas seriam o equivalente brasileiro da currency school e
papelistas o da banking school (p. 697, notas 16 e 17). Itaboraí e Torres Homem
são os expoentes da ortodoxia financeira; Souza Franco, o líder incontestado dos
papelistas (p. 697). A legitimidade intelectual dos dois campos vinha da citação
a autores estrangeiros, com Gilbart, ligado à free banking school e traduzido em
1859, ao lado de Hume, Smith, Ricardo e Mill (p. 698).

\item A demonstração é feita sobre o próprio Souza Franco. Em Os Bancos do
Brasil, de 1848, ele prefere a circulação metálica ou o papel resgatável à vista
às notas inconversíveis do Tesouro, e exige reserva metálica dos bancos (p. 699).
Como ministro, em 1857 e 1858, impôs limites à emissão dos seis bancos que
autorizou. O campeão dos papelistas não estava distante de Torres Homem quanto ao
apoio a uma moeda lastreada em ouro (p. 699).

\chave{The broad metalista vs. papelista dichotomy found in the literature
usually fails to grasp this crucial point in the debate, namely, the defence (or
opposition to) the gold standard}{p. 699}

\item Os rótulos são posteriores aos fatos. Referências contemporâneas às duas
escolas são muito escassas, e os termos teriam sido difundidos por autores do
século XX, com Normano (1935) como exemplo precoce (p. 700, nota 27).

\chave{the terms papelista and metalista apparently having been disseminated by
twentieth-century writers}{p. 700}

A exceção documental recolhida é a definição dada por Itaboraí no Senado em 2 de
junho de 1858, a pedido de Souza Franco: papelistas seriam os que proclamam que o
papel bancário é capital, criam bancos sem impor as condições que garantam o
resgate das notas em ouro e defendem a expansão quando a reserva metálica cai
(p. 700). Mesmo nessa definição hostil, eles continuam partidários de emissão
lastreada em ouro.

\item Os papelistas ``verdadeiros'', defensores de circulação puramente
fiduciária, seriam dois. Mauá, amigo e aliado de Souza Franco, cuja posição
contra a conversibilidade Villela chama de herética para a época (p. 700), e o
estatístico Sebastião Ferreira Soares, para quem nenhum Estado retém ouro em
circulação sem que produção e exportações superem as importações (p. 702). O
autor os aproxima dos economistas de Birmingham, filiação que Jequitinhonha já
usara no Senado como acusação contra Souza Franco (p. 702, nota 35).

\item O desfecho é de derrota parcial da ortodoxia. A Lei dos Entraves de 1860
forçou a contração dos bancos de emissão, o Banco do Brasil retomou a
conversibilidade em outubro de 1862, a crise Souto a suspendeu em setembro de
1864, e em 1866 a Guerra do Paraguai devolveu a emissão ao Tesouro, com retorno
ao padrão-ouro apenas em 1888 (p. 703). Os conservadores financeiros tiveram meia
vitória: não impuseram a circulação integralmente lastreada, objetivo último de
todo metalista, mas restauraram a centralização da emissão (p. 703). A conclusão
retoma Milet, para quem todos os ministros da Fazenda, de Torres Homem a Souza
Franco, eram ``blind admirers of gold'', cegueira partilhada pelos dois campos
(p. 704, nota 37).
\end{enumerate}

\bloco{Conceitos e definições operacionais}
Conversibilidade: resgate das notas em espécie à vista, com o câmbio fixo que ele
implica (p. 693). É o critério classificatório que o artigo propõe em
substituição ao par metalista e papelista. Conversibilidade potencial: promessa
de resgate em ouro ou em papel do Tesouro que, sem lastro metálico integral,
sustenta a credibilidade e mantém o câmbio junto do par (p. 697). Direitos de
emissão contra lastro: os dois eixos independentes do conflito monetário do
século XIX (p. 695). Metalista e papelista: rótulos historiográficos do século
XX, o segundo de uso pejorativo, associado a inflacionistas irresponsáveis
(p. 700). Par legal do mil-réis: 67,5 pence historicamente, 43,4 em 1833, 27 em
1846 (p. 696, nota 15).

\bloco{Evidência e método}
História do pensamento combinada a narrativa institucional, sem modelo nem teste
estatístico. O núcleo temporal vai de 1846 a 1866. As fontes primárias citadas
são relatórios do Ministro da Fazenda de 1850-70, os relatórios das comissões de
inquérito de 1859 e de 1864, os Anais do Senado do Império de 1858, Os Bancos do
Brasil de Souza Franco (1848), o parecer da comissão do Banco do Brasil de 1861,
a autobiografia de Mauá e as Notas Estatísticas de Ferreira Soares (1860). A
Tabela 1 usa cotações diárias da Bolsa do Rio registradas pela Junta dos
Corretores de Fundos Públicos (p. 697). Não há contrafactual.

\bloco{Agentes e vertentes}
Inglaterra: bulionistas e antibulionistas; currency school com McCulloch e Loyd
Overstone; banking school com Mill, Tooke e Fullarton; Gilbart pela free banking
school. Campo metalista brasileiro: Itaboraí e Torres Homem como expoentes, e
entre os políticos Ângelo Muniz da Silva Ferraz, o Visconde de Uruguai e o
Visconde do Rio Branco (p. 697-698, nota 18). Campo papelista: Souza Franco como
líder, com Mauá, Paranaguá e José Antonio Saraiva (nota 18). Contra a
conversibilidade, apenas Mauá e Sebastião Ferreira Soares (p. 700-702). A crítica
nominal de Villela recai sobre Flávio Saes, Crédito e Bancos no Desenvolvimento
da Economia Paulista, e sobre a dissertação de A. M. R. de Andrade sobre 1864
(p. 699, nota 24), esta por propor ``pluralista'' e ``metalista-monopolista''
como substitutos, taxonomia que o autor considera igualmente enganosa (p. 699).

\bloco{Críticas e limites}
A tese é forte e a base documental é estreita. A demonstração se apoia em duas
passagens, o livro de Souza Franco de 1848 e a definição de Itaboraí de 1858, sem
levantamento sistemático de posições nem contagem dos que rejeitavam a
conversibilidade. Se os dois campos convergiam no essencial, fica por explicar a
intensidade do conflito; o autor indica a pluralidade de emissão e as regras
prudenciais como resposta, sem desenvolvê-la. A narrativa para em 1866, com
menção rápida a 1888, e não trata do Encilhamento, da República nem da Caixa de
Conversão: o artigo não sustenta nenhuma afirmação direta sobre 1906. A
afirmação sobre a difusão dos termos por autores do século XX vem com
``apparently'' e um único exemplo (p. 700), e permanece hipótese, não resultado.
Há ainda discrepância de data entre o corpo do texto, que fala em autobiografia
de 1878, e a nota 29, que registra publicação original em 1875.

\begin{dialogo}
\item Vocabulário de época em 1906. Procurar ``metalista'' e ``papelista'' nos
quatro jornais, separando autodesignação de acusação. Villela sustenta que os
termos foram difundidos por autores do século XX (p. 700). Achá-los correntes em
1906 antecipa a difusão e mostra que a clivagem imperial ainda servia como arma
retórica; não achá-los reforça a decisão de tratar o eixo do período como
valorização contra emissão.

\item A conversibilidade como divisor. Verificar se algum participante do debate
sobre a Caixa nega a conversibilidade em si, e não apenas a taxa de 15 pence, o
limite de emissão ou a competência da Caixa. Se ninguém a nega, o debate de 1906
repete a estrutura que Villela descreve para 1858, e ``expansionista'' não pode
ser lido como recusa do padrão-ouro.

\item Conversibilidade potencial e credibilidade. Procurar o argumento de que a
promessa de troco basta para segurar o câmbio mesmo sem lastro integral, e o
contrário, de que só o ouro efetivamente depositado sustenta o mil-réis. Testa se
a formulação de Villela (p. 697) tem correspondente nos debates sobre o fundo de
garantia, o fundo de resgate e a custódia do ouro em 1911.

\item Direitos de emissão contra lastro. Examinar se o conflito entre Caixa,
Tesouro e Banco do Brasil é sobre quem emite, herdeiro da disputa de 1857-1860,
mais do que sobre com que lastro se emite. Se for, o eixo duplo de Villela
(p. 695) reorganiza a leitura do debate de 1906.

\item Linhagem imperial invocada. Procurar menções a Souza Franco, Itaboraí,
Torres Homem, Mauá e ao Encilhamento nos artigos e nos discursos reproduzidos.
Testa se os publicistas de 1906 se filiavam às categorias imperiais ou se a
filiação é construção posterior da historiografia.
\end{dialogo}

\usonocapitulo{Parte III. É a peça que justifica a cautela do capítulo com o par
``ortodoxo'' e ``expansionista'' e com o par de época ``metalista'' e
``papelista''. Sustenta três afirmações. Primeira, que a importação das
categorias inglesas pelo debate imperial foi real e seletiva, porque metalistas e
papelistas herdaram da currency school e da banking school a disputa sobre regra
e discrição sem herdar divergência alguma sobre a conversibilidade, que ambas
pressupunham (p. 693). Segunda, que a clivagem metalista contra papelista, como
usada pela historiografia, não discrimina posições sobre o padrão monetário, e
por isso não pode ser transposta para 1906 como se nomeasse dois programas
opostos. Terceira, que os rótulos são em boa medida invenção classificatória do
século XX (p. 700), o que recomenda tratá-los como vocabulário a verificar nas
fontes, não como categorias analíticas herdadas. O conceito de conversibilidade
potencial (p. 697) alimenta também a Parte IV, ao descrever um arranjo de lastro
parcial que ainda assim produziu estabilidade cambial.}
